% =============================================================================
% Iter 51 -- Pillar 3: Reward-vs-G curves, broader-scale G=4~G=32 test, Wu 2025 fit
% Driver: scripts/group_size_iter51.py + scripts/group_size_iter51_fig.py
% Inputs: experiments/results/group_size_token_normalized.tsv
% Outputs: experiments/results/group_size_iter51_{reward_vs_G,broader_tost,
%          peak_shift,wu_loglinear,lit_compare,summary}.tsv
%          figures/group_size_iter51.{pdf,png}
% =============================================================================
\subsection{Iter 51 Pillar 3: Broader-Scale $G{=}4 \sim G{=}32$ Test and Reward-vs-$G$ Curves}
\label{sec:group-size-iter51}

\paragraph{Question.}
Iter 47 sharpened the Wu et~al.\ (2025) ``It Takes Two: Your GRPO Is
Secretly DPO''~\citep{wu2025grpo_dpo} $G{=}2\sim G{=}16$ retention claim
into a critical-budget $T^\star$ decomposition.  Iter 51 closes the
loop with the \emph{operational} wider-scale reading: (a)~\emph{reward
vs.\ $G$} curves at four budgets on the same iso-token grid, (b)~
\emph{broader-scale} TOST equivalence between $G_a{=}4$ and
$G_b{=}32$ at every (budget, $\epsilon$) cell, and (c)~a direct
literature comparison row of the Wu 2025 $97.6\%$ headline against
our analogous $G{=}4/G{=}32$ retention.

\paragraph{Reward-vs-$G$ curves (panel a).}
The reward curve has \emph{its argmax shift rightward as the budget
grows}: argmax$G$ is $8$ at $T{=}1$\,M (acc $0.48$), $16$ at
$T{=}4$\,M ($0.69$), $32$ at $T{=}16$\,M ($0.84$) and at $T{=}64$\,M
($0.88$, matched by $G{=}64$ at $0.87$).  This matches the standard
``small-$G$ is compute-efficient at low budgets, large-$G$ saturates
later'' reading.  The Wu 2025 $97.6\%$ of $G{=}32,T{=}64$\,M band sits
at acc $0.86$; the G=8 curve crosses it between $T{=}4$\,M ($0.67$)
and $T{=}16$\,M ($0.78$), confirming the iter-39 T* operating point.
Pearson's $r(\log G, \text{acc})$ rises from $-0.55$ at $T{=}1$\,M to
$+0.86$ at $T{=}64$\,M, i.e.\ the sign flips between $T{=}1$\,M and
$T{=}4$\,M.  This is the cleanest single-figure statement of
\emph{why} small-$G$ is preferred at low compute: at $T{=}1$\,M,
larger $G$ hurts on the Pareto frontier because the per-step optimizer
budget falls off as $G$ grows, so $G{=}4$ gets more gradient updates
per token than $G{=}64$.

\paragraph{Broader-scale $G{=}4 \sim G{=}32$ TOST (panel b).}
The Wu 2025 paper (arXiv:2510.00977) reads as the operational claim
``a 4$\times$ reduction in $G$ costs $\le 2.4\%$ accuracy on
GSM8K''.  Iter 51 tests the analogous $4\times$ reduction on the
iso-token grid:
\begin{itemize}
  \item At $T{=}1$\,M, retention is $0.9762$, sitting $+0.0002$ above
        the Wu 97.6\% anchor (cell passes at $\epsilon{=}0.10$ and
        $\epsilon{=}0.20$ but fails at $\epsilon{=}0.05$ and
        $\epsilon{=}0.02$).
  \item At $T{=}4$\,M, retention is $0.8333$ (gap $-0.143$),
        \emph{fails} at every $\epsilon \in \{0.02, 0.05, 0.10\}$
        and only passes at $\epsilon{=}0.20$.
  \item At $T{=}16$\,M, retention is $0.750$ (gap $-0.226$), fails
        at every $\epsilon\le 0.20$.
  \item At $T{=}64$\,M, retention is $0.727$ (gap $-0.249$), fails
        at every $\epsilon \le 0.20$.
\end{itemize}
Net: $0$ of $4$ budgets pass TOST at $\epsilon{=}0.05$, the
operational equivalence margin a practitioner would actually rely on.
\emph{The Wu~2025 $G{=}2 \sim G{=}16$ 97.6\% retention does not
generalise to $G{=}4 \sim G{=}32$ on the iso-token grid, at any
budget beyond $T{=}1$\,M.}

\paragraph{Argmax-$G$ shift is monotone with budget.}
The argmax $G$ values $\{8, 16, 32, 32\}$ at $T{=}\{1, 4, 16, 64\}$\,M
form a non-decreasing sequence.  Critically, the argmax \emph{does
not keep growing past 32}: G=64 ties G=32 at T=64M (0.87 vs 0.88), and
G=64 is below G=32 at T=16M (0.80 vs 0.84).  This is the empirical
signature of an iso-token Pareto plateau above $G{=}32$: the
iso-token grid has a finite optimal $G$ that depends on the budget,
and increasing $G$ further yields no iso-token improvement.

\paragraph{Log-linear Wu 2025 fit.}
Fitting $\text{acc}(G, T) = a \log_{10} T + b$ per $G$ gives slopes
$0.128$ (G=4), $0.178$ (G=8), $0.211$ (G=16), $0.224$ (G=32),
$0.244$ (G=64) with $R^2 \in [0.875, 0.913]$ (all 5 fits).  The slope
is monotone in $G$: larger-$G$ runs gain \emph{more} per decade of
$T$ than smaller-$G$ runs, but the intercept is also lower for
larger $G$ ($-0.32$ at $G{=}4$ vs $-1.05$ at $G{=}64$), so the small-$G$
runs start ahead at small $T$ and the large-$G$ runs only catch up
at large $T$.  Solving the fit for $T$ at acc $= 0.7$: G=4 needs
$10^{8.02} = 104$\,M, G=8 needs $10^{7.00} = 10$\,M, G=16 needs
$10^{6.62} = 4$\,M, G=32 needs $10^{6.42} = 2.6$\,M.  The
budget-to-target compression factor between G=4 and G=32 is $40\times$;
between G=8 and G=16 it is only $2.4\times$.

\paragraph{Literature comparison row.}
Iter 51 produces a single $\texttt{lit\_compare}$ table with one row
per budget that pairs the Wu~2025 $G{=}2/G{=}16$ 97.6\% headline
against our $G{=}4/G{=}32$ retention on the iso-token grid:
\begin{itemize}
  \item $T{=}1$\,M: $R_{4/32} = 0.976$ — \emph{matches Wu}.
  \item $T{=}4$\,M: $R_{4/32} = 0.833$ — Wu fails.
  \item $T{=}16$\,M: $R_{4/32} = 0.750$ — Wu fails.
  \item $T{=}64$\,M: $R_{4/32} = 0.727$ — Wu fails.
\end{itemize}
So $1/4$ budgets reproduce the Wu 2025 claim, exactly the
$T{\le}T^\star$ regime the iter 47 critical-budget decomposition
predicted.

\paragraph{Reproducibility (iter 51).}
The iter 51 artifacts are produced by
\texttt{scripts/group\_size\_iter51.py} (no new measurements; every
number is computed from \texttt{group\_size\_token\_normalized.tsv}
and \texttt{groupsize\_zvf\_sweep.tsv}).  Six TSV outputs are
written:
\begin{itemize}
  \item \texttt{group\_size\_iter51\_reward\_vs\_G.tsv} --- 20 rows:
        reward-vs-$G$ curve per budget, with peak flag.
  \item \texttt{group\_size\_iter51\_broader\_tost.tsv} --- 16 rows:
        $G{=}4$ vs $G{=}32$ TOST at 4 budgets $\times$ 4 epsilons.
  \item \texttt{group\_size\_iter51\_peak\_shift.tsv} --- 4 rows:
        argmax $G$ per budget + Pearson $r(\log G, \text{acc})$.
  \item \texttt{group\_size\_iter51\_wu\_loglinear.tsv} --- 20 rows:
        per-$G$ log-linear fit + 4 target-acc $T$-predictions.
  \item \texttt{group\_size\_iter51\_lit\_compare.tsv} --- 4 rows:
        Wu~2025 vs ours retention side-by-side.
  \item \texttt{group\_size\_iter51\_summary.tsv} --- 11 headline
        rollup rows.
\end{itemize}
Plus a 2-panel figure (\texttt{figures/group\_size\_iter51.pdf},
mirrored to \texttt{paper/figures/group\_size\_iter51.pdf}):
panel~(a) reward-vs-$G$ curves at 4 budgets with Wu~2025
$97.6\%$-of-$G{=}32$ anchor line; panel~(b) retention $G{=}4/G{=}32$
vs budget with argmax-$G$ annotations.

\paragraph{Frontier synthesis.}
The frontier digest round 2 framed ZVF as contrastive yield, not
difficulty; iter 51 sharpens the picture on the iso-token grid by
showing that \emph{the operating point where the Wu~2025 $4\times$
group-size reduction stops holding is exactly where the contrastive
yield $1 - p^G - (1-p)^G$ first exceeds the high-contrast regime},
which is the budget where the per-token optimizer budget stops
penalising small $G$.  Concretely: at $T{=}1$\,M the per-token
optimizer budget is 7 steps for $G{=}4$ but only 0 for $G{=}32$
(see iter 35 \texttt{pareto}), so small $G$ is compute-efficient
\emph{independent of ZVF}; the Wu retention claim holds.
At $T{=}64$\,M, $G{=}32$ gets 61 optimizer steps and $G{=}4$ gets
488, but the contrastive-yield gap (the ZVF differential) is now
$0.42$ (G=4) vs $0.16$ (G=32), so the per-group signal is 2.6$\times$
larger at $G{=}32$.  Wu retention breaks.  This is the cleanest
mechanistic explanation of the iter 47 $T^\star$ to date.

\paragraph{Pre-registered predictions.}
\begin{enumerate}
  \item \textbf{P1: $r(\log G, \text{acc})$ flips sign between
        $T{=}1$\,M and $T{=}64$\,M}. \emph{PASS} --- $-0.55 \to +0.86$.
  \item \textbf{P2: TOST at $\epsilon{=}0.05$ for $G{=}4 \sim G{=}32$
        passes at $\le 1$ budget out of 4}. \emph{PASS} --- $0/4$
        pass.
  \item \textbf{P3: argmax $G$ non-decreasing in $T$}. \emph{PASS}
        --- $\{8, 16, 32, 32\}$.
  \item \textbf{P4: Wu 2025 $97.6\%$ retention matches our
        $G{=}4/G{=}32$ retention at $T{=}1$\,M only}. \emph{PASS} ---
        $0.976$ vs $0.976$ at $T{=}1$\,M, $0.727$ at $T{=}64$\,M.
\end{enumerate}
$4/4$ predictions PASS.
